\section{Introduction}
\label{sec:introduccion}

Ce travail établit une génération coinductive commune de
\(\pi,\varphi,e\), la clôture dodécaphasique de \(\alpha\) et sa
représentation analytique corrélée, ainsi que la composition algébrique de
la coordonnée électronique. Le bloc électronique détermine en outre une
représentation de spin \(1/2\), dont les projecteurs donnent les hélicités
\(\pm1/2\) lorsqu'une direction de propagation est fixée. La même généalogie
conserve un registre d'incidence qui conduit aux systèmes de Witt,
aux codes de Golay, au réseau de Leech et à \(V^\natural\). Le
résultat unit donc des valeurs déterminées, une mémoire récupérable et des
réalisations algébriques par des opérations explicites.

La détermination des valeurs des constantes fondamentales constitue une
question centrale de la physique mathématique. Une théorie peut établir des relations
entre observables, identifier des symétries et organiser des prédictions de grande précision,
tout en conservant des paramètres dont la valeur est déterminée expérimentalement.
Comprendre pourquoi ces paramètres prennent leurs valeurs requiert une construction
qui identifie les opérations, les normalisations et les conditions de
compatibilité dont ils résultent. Telle est la question qui organise le présent
travail : la relation entre la structure qui détermine une constante, l'équation
qui exprime cette détermination et la valeur obtenue par son évaluation.

Cette question est abordée au moyen d'un noyau formel mathématique appelé
\emph{holographie modulaire triadique}, ci-après HMT. Ses trois composantes sont
\emph{All Possible Paths}, ci-après APP ; TRIT, dénomination de la structure
ternaire orientée ; et \emph{Topological Phase Kernel}, ci-après TPK.
APP est le support toroïdal discret \(X_9=(\mathbb Z/9\mathbb Z)^2\).
Ses translations cardinales définissent le graphe de Cayley de la structure
additive, de générateurs \(\pm(1,0),\pm(0,1)\) ; sur ses sommets
sont calculées les évaluations somme et produit, en conservant résidus et
quotients. TRIT est la signature orientée \(\{-1,0,+1\}\), qui distingue
les régimes hyperbolique, parabolique et elliptique. TPK compose la sélection
de phase, le transport des évaluations et la mise à jour de sa
mémoire. Le graphe, les évaluations et la dynamique ont ainsi des fonctions
différentes sur un même état.

La construction détermine la réalisation numérique conjointement avec la structure qui la produit. Les régions de clôture, de propagation et d'auto-échelle sont construites dans un domaine fini de conditions initiales orientées ; leur prolongement compatible détermine ensuite les coordonnées \(\pi,e,\varphi\). Le registre dodécaphasique intervient dans la détermination de \(\alpha\). La construction électronique compose ces coordonnées avec une structure centrale et ses opérateurs de retour. Enfin, l'incidence conservée au cours du parcours admet une réalisation au moyen de codes, de réseaux et d'une algèbre de vertex. Les valeurs et les relations d'incidence sont des réalisations de données communes ; la conservation de ces données pendant le raffinement établit leur connexion mathématique.

La question touche au statut explicatif des constantes. Lorsqu'une
théorie reçoit un paramètre déterminé expérimentalement, ses équations
établissent les conséquences découlant de cette valeur. Une dérivation
structurelle modifie cette relation : la valeur devient l'une des
conséquences de la construction. La question de l'intensité d'un
couplage ou d'un rapport caractéristique de la matière devient
alors celle des relations mathématiques qui les déterminent.
Le passage du paramètre à la conséquence constitue l'argument directeur
de cet article.

La détermination numérique acquiert, dans cette perspective, un contenu
constitutif. L'équation réunit des opérations dont la provenance doit être
explicitée ; son évaluation fournit la valeur que ces opérations fixent. Une
fois établies les conditions correspondantes d'existence, d'unicité et de compatibilité,
la nécessité du résultat relève de la démonstration.
Expliquer pourquoi une constante prend une valeur signifie reconstruire cette
nécessité dans le domaine de la théorie. La précision décimale manifeste
une conséquence quantitative de la structure, tandis que la dérivation
expose la raison pour laquelle cette conséquence précise est obtenue.

La portée de cette question s'accroît lorsque plusieurs déterminations
partagent une même origine. La thèse proposée par les auteurs relie
les constantes à l'organisation produisant leurs relations réciproques.
L'architecture commune doit expliquer tant l'individualité de chaque
coordonnée que les dépendances qui la rattachent aux autres.
L'ordre des constructions acquiert ainsi une valeur démonstrative :
il permet de reconnaître quelles données sont engendrées, quelles opérations les utilisent et
quelles relations subsistent dans les réalisations successives. La diversité
des valeurs est étudiée comme expression d'une unité formelle articulée.

Dans cette organisation, la métrologie remplit une fonction ultérieure et
essentielle. La construction fixe ses résultats avant de les confronter aux
grandeurs observées ; l'expérience examine leur correspondance
physique et délimite la portée de l'accord empirique. La nécessité mathématique s'établit
relativement aux prémisses et aux opérations déclarées, et la réalisation
naturelle est examinée par l'identification des observables et leur
confrontation indépendante. Les deux tâches contribuent à une explication
complète : la première détermine le résultat au sein du corpus ; la seconde
étudie sa pertinence pour décrire la constitution quantitative du
monde.

Cette perspective permet de situer historiquement la recherche. Les
antécédents considérés ci-dessous précisent différentes manières de
relier lois, paramètres et observation. Le présent travail adopte
comme problème central la détermination interne des valeurs et de la
structure qui les relie. Ses développements arithmétiques, géométriques et
algébriques tirent leur unité argumentative de cette question. L'aspiration
ontologique consiste à expliquer comment des propriétés physiques apparemment
indépendantes peuvent correspondre à des relations nécessaires d'une même
organisation mathématique.


Ce problème se présente avec une netteté particulière pour les couplages adimensionnels.
Une coordonnée exprimée dans des unités dépend de la convention métrologique ;
une relation adimensionnelle permet d'interroger sa détermination
indépendamment de ce choix. La constante de structure fine intervient
dans des relations quantitatives susceptibles d'être confrontées à l'expérience et soulève, en même temps, la question
de la structure qui fixe le couplage électromagnétique. Le programme
historique de son explication fournit un antécédent précis pour
étudier conjointement la valeur d'une constante et l'organisation
mathématique dont elle procède.

Un épisode éclairant apparaît dans le travail de Dirac de 1931 sur les
singularités quantifiées du champ électromagnétique. Son objectif concerne
expressément l'existence d'une charge électrique minimale. Le
résultat établit une relation entre cette charge et l'intensité d'un
monopôle magnétique ; Dirac remarque toutefois que la construction ne
détermine pas à elle seule le quotient adimensionnel associé à la valeur
approximative \(137\)\cite{dirac1931}. On obtient une véritable contrainte structurelle
dont la force démonstrative n'équivaut pas à une fixation indépendante
du couplage. Ce précédent montre que la question de l'origine
des constantes appartenait à la recherche fondamentale et exige
d'identifier exactement ce que détermine chaque théorie.

L'histoire de l'électrodynamique quantique précise la différence entre
cette exigence et le calcul d'effets observables. Dans sa conférence Nobel
de 1965, Feynman a reconstitué des procédures exprimant les déplacements
de niveaux et les processus de diffusion au moyen de la masse expérimentale de
l'électron\cite{feynman1965}. La renormalisation a permis d'obtenir
des résultats finis et confrontables à l'expérience tout en conservant cette référence. L'efficacité
de ces prédictions constitue un acquis théorique effectif ; leur dépendance
à l'égard d'un paramètre expérimental identifie, en même temps, la place logique
où une détermination supplémentaire pourrait élargir l'explication. La
question constitutive s'ajoute à la capacité prédictive et étudie les
conditions qui fixent ses paramètres.

Le rapport conjoint LEP--SLD de 2006 fournit une manifestation expérimentale
de cette architecture\cite{lep2006}. Ses auteurs distinguent les paramètres
libres du modèle des relations que celui-ci impose entre observables.
Par les corrections radiatives, les données de la résonance Z permettaient
d'inférer indirectement les masses du quark top et du boson W et de les confronter
aux mesures directes. Les prédictions entrelacées et la détermination
expérimentale des paramètres formaient une même procédure de mise à l'épreuve.
Le problème pertinent consiste donc à identifier les dépendances qu'une
construction détermine intérieurement et celles dont la spécification continue
de nécessiter une information empirique.

Dans ce contexte, on entendra par \emph{détermination épistémologique}
la procédure qui établit une valeur par la relation entre théorie,
observation et inférence ; par \emph{détermination constitutive}, la
déduction de cette valeur à partir des relations définissant l'objet
considéré. Cette seconde expression désigne ici une exigence mathématique :
expliciter le domaine, les opérateurs, les normalisations et les
compatibilités sélectionnant le résultat. Sa pertinence physique
procède de la confrontation ultérieure aux grandeurs observées.
Le changement de perspective concerne l'ordre de l'explication : la valeur est
examinée comme conséquence d'une structure déterminée. L'aspiration
ontologique consiste à relier cette structure à la constitution de
l'objet physique ; la preuve mathématique et sa confrontation expérimentale établissent
la portée avec laquelle une telle relation peut être soutenue.



La constante de structure fine rend particulièrement manifeste la portée du
problème. Dans la description électromagnétique usuelle, \(\alpha\) est un
couplage sans dimension : elle exprime une relation physique qui subsiste lors
d'un changement d'unités. Pauli associait l'explication de sa valeur à la
compréhension de la structure atomistique de l'électricité\cite{pauli1946}.
On présente ici d'abord la construction mathématique du couplage puis
celle de l'électron auquel cette interaction est attribuée. Leur proximité
dans l'argument répond à une dépendance structurelle : les données qui
déterminent le couplage interviennent également dans la composition de la
coordonnée électronique. La comparaison physique examine en dernier lieu ces
expressions préalablement fixées.

La perspective historique comprend également la relation entre action, phase
et orientation. Dans sa formulation par chemins, Feynman attribue aux
trajectoires une phase dépendant de l'action et obtient l'amplitude par
composition de leurs contributions\cite{feynman1948}. La dénomination
\emph{All Possible Paths} reprend cette référence conceptuelle. Dans la construction
présente, les chemins sont définis sur une structure arithmétique finie ;
les règles d'admissibilité et leurs évaluations sont spécifiées par
des opérations entières. La relation avec une action physique est formulée à
l'étape dimensionnelle, une fois ce domaine construit.

Le lien entre action et phase angulaire apparaît également dans le traitement
lagrangien de Dirac\cite{dirac1933}. L'expression \(\hbar=h/(2\pi)\)
situe l'action par rapport à une coordonnée angulaire en radians. Le développement
actuel étudie préalablement les déterminants, quotients et transports qui
produisent une échelle d'action sans dimension. Son application à une unité
physique se distingue des rapports dans lesquels cette échelle commune se simplifie.
Cette distinction est décisive pour expliquer la composition électronique
et le rôle du retour chronologique : le quantum d'horloge fournit une
échelle, tandis que la section électronique conserve son identité structurelle.

Dans sa conférence Nobel de 1965, Feynman rappelait une conversation avec
John Archibald Wheeler au sujet de l'identité de charge et de masse des
électrons\cite{feynman1965}. Wheeler imaginait une unique ligne d'univers
qui, parcourant le temps selon des orientations alternées, apparaîtrait dans une
section temporelle comme de multiples électrons et positrons. Feynman distinguait
cette conjecture de l'électron unique de la représentation qu'il incorpora à
son travail : un positron comme segment de ligne électronique d'orientation
temporelle inverse. Son article de 1949 développe cette représentation et
reconnaît l'antécédent de Stückelberg\cite{feynman1949}. Cet épisode
situe historiquement le problème de l'orientation des trajectoires. En HMT,
la bidirectionnalité s'exprime par des involutions, des transports et
des conditions de retour dont le domaine est spécifié dans chaque construction.

Le contenu arithmétique commence par la coexistence de deux évaluations.
Chaque ligne de la table multiplicative s'obtient par accumulation additive ;
la réduction modulaire enregistre ses résidus, tandis que les quotients conservent
les multiples de la base qui franchissent la frontière. La différence entre
somme et produit acquiert ainsi une expression exacte sur un même support.
Les compositions mixtes peuvent être sensibles à l'ordre et engendrer une
courbure discrète. TRIT distingue les régimes elliptique, parabolique et
hyperbolique par leurs relations quadratiques, et l'exponentielle fournit
une loi commune aux trois. Cette organisation permet d'étudier un nombre
comme réalisation scalaire d'une section compatible conservant davantage de données
que sa seule coordonnée réelle.

Le raffinement explicite cette distinction. Une suite compatible de
cylindres de diamètre décroissant détermine une coordonnée archimédienne ;
la trajectoire conserve simultanément orientation, quotients de transport,
régime et mémoire. L'holonomie nonadique identifie le retour de phase,
tandis que son relèvement enregistre l'accroissement de profondeur. La représentation
sturmienne décrit l'ordre apériodique de ce prolongement et fournit
une formulation opératorielle de ses deux phases. La portée de cette construction
symbolique est établie par ses opérateurs et ses invariants.

Le registre \(K\) reçoit une définition propre avant d'intervenir dans
\(\alpha\). Il est construit par une transformation entière réversible du
registre signé ; sa réalisation rationnelle \(\kappa\) conserve les douze
composantes, l'origine et l'orientation fixées. Les différences cycliques
déterminent les relations sectorielles, et la composante uniforme complète
sa reconstruction. Pour \(\alpha\), on distingue la détermination positionnelle
par normalisation entière et la détermination analytique par le résolvant
des lacunes. Leur identification exige de conserver la compatibilité des
prolongements. Les conséquences relatives à \(e+\pi\) seront exprimées en termes
de ces relations exactes et des conditions arithmétiques qu'elles permettent
effectivement de décider.

La sélection unique de \(K\) est établie par des opérations sur
des populations explicites. La suite
\(40320\to21902\to19446\to79\to1\) passe des ordonnancements aux
incidences indexées avec la grille, aux numérateurs distincts, aux candidats de
la face exceptionnelle et aux survivants hétérotypiques. Le recensement régional
\(196\cdot197\cdot196\to331\to2\to1\) applique des conditions de Gram,
d'occupation, de distance et d'orientation à des blocs de trois catalogues.
La section~\ref{act21:sec:relacion-censos} distingue les domaines des
deux dénombrements ; les démonstrations de sélection et d'inversion précisent
comment leurs sorties se composent avec le registre réversible.

La section électronique développe l'origine de chaque facteur de l'équation :
la norme matricielle \(\sqrt3/4\), le transfert régional
\(\varphi/\pi^2\), le terme cubique en \(\alpha\), la normalisation
torsionnelle et la mémoire multiplicative de retour. Deux lecteurs de la
trajectoire centrale déterminent le même triplet de coefficients par
des opérations distinctes. La même récurrence détermine un support de
\(27\) cellules et deux historiques orientés que la coordinatisation intégrale distingue.
Sa lecture harmonique identifie trois modes du registre central et exprime
les coefficients du retour au moyen de leurs puissances. L'équation électronique
conserve ainsi une dépendance explicite à l'égard de la trajectoire produite.
La structure locale et la coordonnée complète sont
maintenues lors de la réalisation dimensionnelle. La comparaison est présentée
avec les ajustements CODATA 2018 et 2022\cite{codata2018,codata2022}, en indiquant
les valeurs centrales, leurs incertitudes et les différences calculées.

Le prolongement exceptionnel reprend les données de frontière de ce même
parcours. Les changements de représentation passent des séquences ternaires aux
codes, des supports aux plans d'incidence et des classes discriminantes aux
réseaux. La construction du voisin de Leech permet d'appliquer le
procédé de Frenkel, Lepowsky et Meurman pour obtenir \(V^\natural\) ;
le Monstre agit comme groupe d'automorphismes de cette algèbre, et Moonshine
décrit ses traces graduées\cite{flm,borcherds}. La réalisation scalaire
et la réalisation d'incidence constituent les deux projections examinées
dans l'article : toutes deux partent de données identifiées et conservent les
relations pertinentes de la construction.

% BEGIN S27_FRAMING
Deux développements précisent l'articulation entre les registres entiers
et leurs incidences. Le paragraphe~\ref{carry27:seccion} dérive le cocycle
de retenue de la division euclidienne des colonnes régionales. Le tableau
complet de \(27\) entrées, les représentations polynomiales et les
preuves de rang établissent la contribution algébrique des quotients.
Les identités de transport conservent cette contribution avec le
résidu ; les bornes du cylindre déterminent la compatibilité de la
lecture décimale.

La section~\ref{star27:comparacion} compare ensuite les étoiles de
\(B_K=H_e\cap P_\pi\), \(B_\varphi=H_\varphi\cap P_\pi\) et
\(B_{\rm APP}=H_{\rm APP}\cap P_\pi\), polarisées par le support
négatif de la cochaîne antérieure à la retenue. Leurs douze complétions,
le recensement des \(495\) quadruplets et les stabilisateurs du cadre
ordonné spécifient l'information retenue par cette lecture
d'incidence. La connexion entre les deux développements réside dans les
registres régionaux et orientés produits par APP--TRIT--TPK :
l'évaluation entière conserve les grandeurs et les retenues, et la réalisation
de frontière conserve les supports, les intersections et les symétries.

% END S27_FRAMING


% BEGIN R27_FRAMING_INTRO_ES
La construction finie permet de préciser quelle information se transmet à
chacune de ces étapes. L'alphabet d'une fenêtre comprend \(42\)
blocs, obtenus à partir des deux signatures de l'émetteur ; les six fenêtres produisent
le catalogue sur lequel s'effectuent les raccordements régionaux. L'exemple
de deux registres ayant le même histogramme et des corrélations temporelles distinctes
identifie une contribution de l'ordre que la lecture marginale omet. Les
quatre jauges de l'observable multiplicatif et les contrôles de clôture
sont examinés sur des domaines fixés, en conservant la provenance du registre
et ses conditions de frontière.

La même discipline distingue trois extensions d'APP : variation du
module, répétition périodique du support et augmentation du nombre de
coordonnées. La compression multiplicative en trois classes conserve un
sous-réseau d'indice deux dont la composante de rang deux réalise une
itération de Pell. Le codage de paires de trits transporte la somme
par sa retenue. Ensuite, la corrélation quadratique construit
l'identité de la coordinatisation à six positions ; la neutralité duale réduit
à deux les huit frontières saturées et l'orientation fixée sélectionne son
représentant. Le résultat de divisibilité de niveau trois se situe au
terme de la réalisation exceptionnelle, où sa démonstration utilise
l'identité rationnelle déjà établie.
% END R27_FRAMING_INTRO_ES

% BEGIN R28_FRAMING_INTRO_ES
La liaison régionale est en outre développée sur le multigraphe étiqueté de l'émetteur et son quotient de phase. Cinq parcours de longueur minimale huit réunissent, dans chaque région de clôture, les ancrages APP, de propagation et d'auto-échelle. Leurs témoins orientés et leur preuve finie de minimalité complètent la description des raccordements et de l'information conservée par chaque représentation.
% END R28_FRAMING_INTRO_ES


L'énumération des conditions initiales délimite les états admissibles ;
les certificats de calcul permettent de vérifier leurs relations finies et
les registres de provenance identifient les données de chaque opération.
Les définitions, les lemmes et leurs démonstrations établissent la continuité
entre ces objets. Les passages à la limite conservent leurs hypothèses et leurs
applications de restriction.
La construction de la section~\ref{subsec:alpha-publicaciones-completas}
spécifie tous les coefficients de la représentation analytique corrélée
et démontre sa compatibilité avec la clôture entière. L'égalité des
préfixes des deux représentations à toute profondeur est établie dans le
théorème~\ref{thm:alpha-dos-publicaciones-completas}, au sein de la classe
géométrique et du domaine qui y sont définis.
La question directrice demeure tout au long
du parcours : quelle structure détermine la valeur et quelles relations subsistent
lorsqu'on change sa représentation.
